\section{Espaces vectoriels, dimensions et applications linéaires}
\subsection{Exercices méthodes :}
\subsection{Exercices classiques :}
\subsubsection{Indépendance du choix de base.}

Soit $B=(x_1,\dots,x_n)$  une base de $E$.
et $A=(a_1,\dots,a_n)$ une autre base.
On note $det_{A}(u) = \lambda_{A}$

$$
\begin{aligned} 
\lambda_A &= det_A(u(a_1),\dots,u(a_n)) \\
\text{par changement de base }&= det_A(B)det_B(u(a_1),\dots,u(a_n))
\\ &= det_A(B)\lambda_B det_B(a_1,\dots,a_n)
\\ &= \underbrace{det_A(B)det_B(A)}_{=1}\lambda_B 
\\ &= \lambda_B 
\end{aligned}
$$

\subsubsection{Cours : $det(A) = det(A^T)$}

On a $$ 
\begin{aligned}
    det(A) &= \sum_{\sigma \in \mathfrak{S}_n} \epsilon(\sigma) \prod_{i=1}^{n} (A)_{i,\sigma(i)} 
    \\ &= \sum_{\sigma \in \mathfrak{S}_n} \epsilon(\sigma) \prod_{i=1}^{n} (A)_{\sigma(i),i} 
    \\ &= \sum_{\sigma \in \mathfrak{S}_n} \epsilon(\sigma) \prod_{i=1}^{n} (A^T)_{i,\sigma(i)} 
    \\ &= det(A^T) \hspace{10px } \square
\end{aligned}
$$